% ---------------------------------------------------------------------
%                      DOCUMENTO DE EJEMPLO
%                    QUE UTILIZA LA PLANTILLA
%                      DE LA EDITORIAL DE LA
%                UNIVERSITAT POLITÈCNICA DE VALÈNCIA
% ---------------------------------------------------------------------

\documentclass[tesi, crop, rm, english]{configuraciones/editorialupv}

% Opciones de la clase 'editorialupv'
%
% llibre      - Libro docente
% ebookpdf    - Libro en formato PDF para e-Readers
% tesi        - Formato de tesis
% a4          - Formato A4 para imprimir
% a4e         - Formato A4 electrónico (márgenes simétricos)
% rm          - Tipo de letra roman
% sf          - Tipo de letra sanserif
% crop        - Marcas de corte (cruz) 17x24 cm (+3 mm por los cuatro lados) centrada en una A4 
% nocrop      - Sin marcas de corte, respeta el tamaño 17x24 cm 
% nomathskip  - No se modifican las distancias de las ecuaciones 
%
% castellano  - La publicación está escrita en castellano
% valencia    - La publicació està escrita en valencià
% english     - La publicación está escrita en inglés
%
% literari    - La publicación tiene un formato literario

% ---------------------------------------------------------------------
% Configuraciones presonalizadas 

% ==============================================
%  PREAMBULO (THESIS SPECIFIC LOGIC)
% ==============================================

% 1. Handle Language Logic (Relies on editorialupv.cls flags)
\ifEPUB
    \usepackage[english]{babel}
\else
    \usepackage[english]{babel}
\fi

% 2. Load the Safe Common Packages
% ==============================================
%  PAQUETES COMUNES (SAFE FOR ARTICLES & THESIS)
% ==============================================

% --- ENCODING & FONTS ---
\usepackage[utf8]{inputenc}
\usepackage[T1]{fontenc}
\usepackage{lmodern}
\usepackage{csquotes}

% --- MATH & SCIENCE ---
\usepackage{amsmath, amssymb, mathtools}
\usepackage{siunitx}
\sisetup{output-decimal-marker={.}}

% --- TABLES ---
\usepackage{array, booktabs, tabularx, longtable, multicol}

% --- IMAGES & COLORS ---
\usepackage{xcolor}
\usepackage{float}
\usepackage{caption}
\usepackage{subcaption}
\usepackage{graphicx} 
\newcommand{\incluyeGrafico}[2][]{\includegraphics[#1]{#2}}

% --- TIKZ & PLOTS ---
\usepackage{tikz}
\usepackage{pgfplots}
\pgfplotsset{compat=newest}
\usetikzlibrary{positioning, shapes, arrows.meta, calc, chains}

% --- UTILITIES ---
\usepackage{lipsum}
\usepackage{import}
\usepackage{collect}
% --- BIBLIOGRAPHY ---
\usepackage[
    backend=biber,
    style=apa, % Change to numeric or ieee if needed for articles
    hyperref=true,
    backref=true,
    url=false
]{biblatex}

% --- COLORS ---
\definecolor{colorEnlace}{rgb}{0,0,0.6}


% Other imports, specific to the articles


\subimport{3DFulltruckRTI/}{necessarypackages}



% \input{./configuraciones/preambulo_tcolorbox}
% \input{./configuraciones/preambulo_listings_LaTeX} % Para incluir código LaTeX
%\input{preambulo_listings} % Ejemplos de configuración para incluir código C, o Matlab

\usepackage{import}
\usepackage{rotating}
\usetikzlibrary{positioning,chains}

\usepackage{tikz}
\usepackage{longtable}
\usepackage{booktabs}
\usepackage{amsmath}
\usepackage{graphicx}
\usepackage{placeins}
\usepackage{tikz}
\usepackage{algorithm}
\usepackage{algpseudocode}
\usetikzlibrary{positioning}
\usepackage{adjustbox}
\usepackage[flushleft]{threeparttable}
\usepackage[singlelinecheck=false, labelfont=bf]{caption}
\usepackage{array}

% ¿Quieres cambiar el color de los enlaces?
% Entonces descomenta la siguiente línea de código y configura el color

% \colorlet{colorEnlace}{blue} 

% ---------------------------------------------------------------------
% Archivos con las referencias bibliográficas


% ---------------------------------------------------------------------
% Las carpetas para los gráficos

\graphicspath{
	{./figuras/}
	{./logos/}
	}

% ---------------------------------------------------------------------

\title{
		Algoritmos de optimización de problemas de Producción Industrial y Logística en la Industria 4.0}


\author{
	\parbox{\textwidth}{\centering
		\begin{tabular}{l@{\hspace{.75em}}l}
			Autor: 		& Juan Moreno Malo\\[4ex]
			Director: 	& D. Josefa Mula Bru\\[.5ex]
						& D. Raul Poler
		\end{tabular}
		}
	}

% ---------------------------------------------------------------------
% ---------------------------------------------------------------------
% ---------------------------------------------------------------------
\addbibresource{3DFulltruckRTI/references.bib}
\addbibresource{RouteClustering/references.bib}




\begin{document}
	
\usepackage{hyperref}
\hypersetup{
    colorlinks=true,
    linkcolor=colorEnlace,
    citecolor=colorEnlace,
    urlcolor=colorEnlace,
    bookmarksnumbered=true,
    breaklinks=true,
    pdftitle={Thesis}, % Fixed string to avoid \abstractname issues
    pdfauthor={Juan Moreno}
}


% -------------------------------------------------------

\frontmatter

% -------------------------------------------------------
% Página de título

\maketitle

% -------------------------------------------------------
% Resumen, prólogo o prefacio

\cleardoublepage


	
% -------------------------------------------------------
% Índice: tabla de contenidos

% Para seleccionar el nivel de títulos en la tabla de contenidos
\setcounter{tocdepth}{1}
% 1: Hasta \section{}
% 2: Hasta \subsection{}
% 3: Hasta \subsubsection{} (No recomendable)    

\cleardoublepage
\phantomsection
\addcontentsline{toc}{chapter}{\contentsname}
\tableofcontents
        
% -------------------------------------------------------
% Lista de figuras
\cleardoublepage
\phantomsection
\addcontentsline{toc}{chapter}{\listfigurename}
\listoffigures

% -------------------------------------------------------
% Lista de tablas
\cleardoublepage
\phantomsection
\addcontentsline{toc}{chapter}{\listtablename}
\listoftables

\cleardoublepage

% -------------------------------------------------------
% Numeración de páginas números arábicos
% Primer capítulo en página 1

\mainmatter

% -------------------------------------------------------
% -------------------------------------------------------
% -------------------------------------------------------
% Capítulos de la publicación

\subimport{RouteClustering/}{RouteClustering.tex}
\subimport{LiteratureReviewRTIs/}{main.tex}

\chapter{An optimisation approach for supply chain procurement transport operational planning with 3D full-truck loading constraints and returnable transport items}
\subimport{3DFulltruckRTI/}{main.tex}


% -------------------------------------------------------
% -------------------------------------------------------
% -------------------------------------------------------
% Bibliografía


% -------------------------------------------------------
% Índice alfabético

\cleardoublepage
\phantomsection
\addcontentsline{toc}{chapter}{\indexname}

\printindex


% -------------------------------------------------------
% Fin del documento

\end{document}

% ---------------------------------------------------------------------
% ---------------------------------------------------------------------
% ---------------------------------------------------------------------
